\section{The abstract elementary class}\label{sec:aec}

We now turn the tests of Section~\ref{sec:colors} into a class of
structures.  The distinction between the object condition and the strong
substructure condition is essential: objects allow countably many bad
coordinates for each tuple, whereas a strong extension must introduce no
bad coordinate for an old tuple outside the old base order.

\subsection{The axioms and the basic structures}

\begin{definition}[Abstract elementary class]\label{aec:definition}
An \emph{abstract elementary class} is a pair
\(K=(\mathcal C,\le_K)\), where \(\mathcal C\) is a class of set-sized
structures in one fixed set-sized finitary language \(\tau\), and
\(\le_K\) is a partial order on \(\mathcal C\), with the following
properties.  Write \(|M|\) for a structure's underlying set and
\(\|M\|\) for its cardinality.
\begin{enumerate}[label=\textup{(\arabic*)}]
\item \emph{Isomorphisms and substructures.}
The class is closed under \(\tau\)-isomorphisms; \(M\le_K N\) implies
that \(M\) is a \(\tau\)-substructure of \(N\); and if
\(f:N\cong N'\) and \(M\le_K N\), then \(f[M]\le_K N'\).
\item \emph{Coherence.}
If \(M_0\) is a \(\tau\)-substructure of \(M_1\), and
\(M_0\le_K M_2\) and \(M_1\le_K M_2\), then \(M_0\le_K M_1\).
\item \emph{Unions and smoothness.}
For every nonzero ordinal \(\gamma\) and every increasing chain
\((M_i:i<\gamma)\) under \(\le_K\), its union \(U\) belongs to
\(\mathcal C\), and \(M_i\le_K U\) for all \(i<\gamma\).
If also \(M_i\le_K N\) for every \(i<\gamma\), where
\(N\in\mathcal C\), then \(U\le_K N\).
\item \emph{Downward L\"owenheim--Skolem property.}
There is an infinite cardinal
\(\theta\ge|\tau|+\aleph_0\) such that, for every \(M\in\mathcal C\)
and every \(A\subseteq |M|\), there exists \(N\le_K M\) with
\(A\subseteq |N|\) and \(\|N\|\le |A|+\theta\).
\end{enumerate}
The least such infinite cardinal \(\theta\) is denoted by \(\LS(K)\).
\end{definition}

We will prove the union and common-upper-bound assertions directly for
every nonempty set-indexed directed family of strong substructures.  This
in particular proves all the chain assertions in
Definition~\ref{aec:definition}, with no restriction on the chain's
cofinality.

For a nonempty set \(I\), let \(\mathcal T(I)\) be the set of all
nonempty finite ordered tuples of distinct members of \(I\).
If \(t=(a_1,\ldots,a_k)\) and
\(\varnothing\ne J\subseteq[k]\), write \(t\restriction J\) for the subsequence in
the original order.  No order is placed on the sort \(I\).

Finite-support groups and torsors also occur in the Hart--Shelah
construction \cite[Section~1]{HartShelah1990};
see also \cite[Definition~3.1 and Fact~3.2]{BoneyVasey2018HartShelah}.
We specify two families of torsors, one for rows and one for columns,
and the precise finite-support action on their common bit relation.

\begin{definition}[Basic structures]\label{aec:basic}
A basic structure has the following data, with pairwise disjoint
underlying carriers.
\begin{enumerate}[label=\textup{(\alph*)}]
\item A nonempty set \(I\), and a nonempty dense linear order \(B\)
without endpoints in which every bounded interval is countable.
\item A group sort \(G(B)\) coding \([B]^{<\omega}\) under symmetric
difference.  More precisely, an incidence relation determines a
bijection
\[
 \supp:G(B)\longrightarrow[B]^{<\omega},
 \qquad
 \supp(u+u')=\supp(u)\SymDiff\supp(u').
\]
There is also a group sort \(G_*\) identified, by a separate unary
singleton label for each of its elements, with
\([\omega]^{<\omega}\) under symmetric difference.  These labels
exhaust \(G_*\).
\item For each \(t\in\mathcal T(I)\) and \(n<\omega\), a nonempty
set \(Z_{t,n}\) with a free transitive action of \(G(B)\).  For each
\(t\in\mathcal T(I)\) and \(b\in B\), a nonempty set \(W_{t,b}\)
with a free transitive action of \(G_*\).  We write both actions
additively.  Thus each element of a fiber is obtained uniquely by
translating any one chosen element of that fiber.
\item A binary relation on row and column points, written as its
indicator \(e(z,w)\in\{0,1\}\).  It can hold only when the row and
column have the same tuple \(t\).  On these pairs it satisfies
\begin{equation}\label{aec:bit}
 e(z+u,w+v)
 =
 e(z,w)+\boldsymbol 1_{b\in\supp(u)}
       +\boldsymbol 1_{n\in v}\pmod 2
\end{equation}
for \(z\in Z_{t,n}\), \(w\in W_{t,b}\), \(u\in G(B)\), and
\(v\in G_*\).  In the last term \(v\) is read through its fixed
finite-subset label. Here \(\boldsymbol 1_P\) is \(1\) when
\(P\) holds and \(0\) otherwise.
\end{enumerate}
The domain consists exactly of the listed carriers.  There are no
additional relations or distinguished representatives.
\end{definition}

Here is an explicit countable finitary relational language for these
data.  Use unary predicates for \(I,B,G(B),G_*\), for each carrier
\[
 Z_{k,n}=\bigcup_{\substack{t\in\mathcal T(I)\\|t|=k}}Z_{t,n}
 \quad(k\ge1,\ n<\omega),
 \qquad
 W_k=\bigcup_{\substack{t\in\mathcal T(I)\\|t|=k,\ b\in B}}W_{t,b},
\]
and for each singleton label in \(G_*\).  The carrier predicates,
excluding the labels within \(G_*\), partition the domain.
For each \(k,n\), a \((k+1)\)-ary projection graph associates a
row point with its \(k\) tuple entries.  For each \(k\), a
\((k+2)\)-ary projection graph associates a column point with its
tuple entries and its coordinate \(b\).  These projections are
semantically total and single-valued on their carriers, have distinct
tuple coordinates in \(I\), and have exactly the nonempty fibers
required in Definition~\ref{aec:basic}.  A carrier may be empty if
there are no tuples of its length.

Use binary relations for the order, finite-set incidence, and the
bit relation, and ternary graphs for group operations and the two
actions.  Every operation and action is semantically total on its
stated domain.  Operation graphs give precisely the specified groups;
action graphs preserve the indicated fibers and give precisely the
free transitive actions.  There are no relation instances outside
these specified types.  In particular, bits between different tuple
fibers are zero.  This last requirement rules out any unspecified
structure linking different tuples.

All the symbols just listed are indexed by countable sets, and each
has finite arity.  The language is therefore countable and finitary.
Using operation graphs is harmless: in an inclusion of induced
relational substructures, a graph triple already in the smaller
object remains in the larger one.  Semantic uniqueness in the larger
object forces it to give the same operation or action value on the
old arguments.  The smaller object's semantic totality supplies the
needed closure on those arguments.

The fixed coding maps and tests from Section~\ref{sec:colors} are set
parameters in the definition of the class, not additional language
symbols.  Neither the countability restrictions nor the condition
using these tests needs to be first-order expressible.

\begin{definition}[Basic inclusion]\label{aec:basic-inclusion}
For basic structures \(M_0,M_1\), write
\(M_0\preccurlyeq_{\mathrm b}M_1\) if \(M_0\) is an induced
\(\tau\)-substructure of \(M_1\), and the following further conditions
hold:
\begin{enumerate}[label=\textup{(\roman*)}]
\item \(G_*^{M_0}=G_*^{M_1}\) as underlying sets;
\item \(B^{M_0}\) is a convex suborder of \(B^{M_1}\);
\item for every \(u\in G(B)^{M_0}\),
\(\supp^{M_1}(u)=\supp^{M_0}(u)\).
\end{enumerate}
All inclusions here and below allow equality.
\end{definition}

Condition (iii) forbids an old finite-set code from acquiring new
members in the larger base order.  Since support gives a bijection
onto finite subsets in both objects, this condition makes the old
group exactly the naturally included copy of
\([B^{M_0}]^{<\omega}\) in \([B^{M_1}]^{<\omega}\).  The group
elements themselves may have arbitrary underlying codes.
Basic inclusions compose: induced substructures, convexity,
equality of \(G_*\), and support preservation each compose.

\begin{lemma}[Size of the base order]\label{aec:order-size}
Every order \(B\) in a basic structure has cardinality at most
\(\aleph_1\).
\end{lemma}

\begin{proof}
Fix \(p\in B\).  If the final segment above \(p\) has a countable
cofinal subset, its cardinality is countable: it is covered by
countably many bounded intervals with one endpoint \(p\), together
with their endpoints.  Otherwise one can recursively choose a
strictly increasing sequence
\((b_\alpha:\alpha<\omega_1)\) above \(p\).  At each stage there
are only countably many earlier choices, so the assumption that
no countable set is cofinal provides a point above all of them.
This sequence must be cofinal above \(p\).  A point above the
whole sequence would enclose uncountably many \(b_\alpha\)'s
in a bounded interval, contrary to local countability.
The final segment is therefore covered by at most \(\aleph_1\)
countable bounded intervals and their endpoints.  It has size at
most \(\aleph_1\).  The same argument in the reversed order
applies below \(p\), proving the claim.
\end{proof}

\begin{lemma}[Size of a basic structure]\label{aec:size}
For every basic structure \(M\),
\begin{equation}\label{aec:size-equation}
 \|M\|=\max\bigl(|I^M|,|B^M|,\aleph_0\bigr).
\end{equation}
\end{lemma}

\begin{proof}
The right-hand side is a lower bound, since \(I,B,G_*\) are
sorts and \(G_*\) is countably infinite.  For the upper bound put
\(\kappa=\max(|I|,|B|,\aleph_0)\).
There are at most \(\kappa\) finite injective tuples over \(I\).
The order \(B\) is infinite, so \(|G(B)|=|B|\); hence every
row fiber has size \(|B|\), and every column fiber has size
\(|G_*|=\aleph_0\).
There are at most \(\kappa\) row fibers, indexed by
\(\mathcal T(I)\times\omega\), and at most \(\kappa\) column
fibers, indexed by \(\mathcal T(I)\times B\).
All the carriers together consequently have size at most
\(\kappa\).
\end{proof}

\subsection{Colors, objects, and strong inclusions}

A \emph{row-choice system} in a basic structure is a family
\(\mathbf z=(z_{t,n}:t\in\mathcal T(I),\,n<\omega)\) with
\(z_{t,n}\in Z_{t,n}\).  Such systems exist by choice, since the
structure and the indexing set are set-sized.  For \(b\in B\) define
\begin{equation}\label{aec:color}
 c_t^{\mathbf z}(b)
 =
 \bigl[(e(z_{t,n},w))_{n<\omega}\bigr]_{=^*}\in S,
 \qquad w\in W_{t,b}.
\end{equation}

\begin{lemma}[Representative independence]\label{aec:representatives}
The color in \eqref{aec:color} is independent of \(w\).
If \(\mathbf z,\mathbf z'\) are two row-choice systems, then,
for each fixed \(t\), their colors \(c_t^{\mathbf z}(b)\) and
\(c_t^{\mathbf z'}(b)\) agree outside a countable subset of \(B\).
The full subsequence test for a fixed tuple therefore has the
same truth value for the two systems outside a countable set.

More strongly, suppose \(M_0\preccurlyeq_{\mathrm b}M_1\).
For a tuple in \(\mathcal T(I^{M_0})\) and a point
\(b\in B^{M_1}\setminus B^{M_0}\), compute colors in \(M_1\)
using row representatives from \(M_0\).  These colors are
independent of every such row choice, with no exceptional
coordinates.
\end{lemma}

\begin{proof}
Two column representatives differ by a translator \(v\in G_*\).
Equation~\eqref{aec:bit} changes their binary sequences only at
the finitely many indices in \(v\), so their classes modulo
\(=^*\) agree.

For a fixed tuple \(t\), write
\(z'_{t,n}=z_{t,n}+u_{t,n}\), with
\(u_{t,n}\in G(B)\).  Put
\[
 D_t=\bigcup_{n<\omega}\supp(u_{t,n}).
\]
Each support is finite, so \(D_t\) is countable.  At every
\(b\notin D_t\), all coordinates of the two bit sequences are
identical when computed with the same column representative.
The full test for a tuple of length \(k\) involves just its
\(2^k-1\) nonempty subsequences.  Outside the union of their
sets \(D_{t\restriction J}\), all its input colors agree, proving the
test assertion.

In the last assertion, the translators between two old
representatives belong to \(G(B)^{M_0}\).  Their supports in
\(M_1\) are still subsets of \(B^{M_0}\), by basic inclusion.
At any \(b\) outside \(B^{M_0}\), Equation~\eqref{aec:bit}
therefore leaves every individual bit unchanged.  This proves
the stronger assertion.
\end{proof}

For a row-choice system \(\mathbf z\) and a tuple
\(t=(a_1,\ldots,a_k)\), let
\begin{equation}\label{aec:bad-set}
 E_t(\mathbf z)=
 \left\{b\in B:
 \neg T_k\bigl((c_{t\restriction J}^{\mathbf z}(b))_{\varnothing\ne J\subseteq[k]}\bigr)
 \right\},
\end{equation}
where \(T_k\) is the full subsequence test of
Definition~\ref{col:test}.

\begin{definition}[Objects and strong substructures]\label{aec:objects}
A basic structure is an \emph{object} if \(E_t(\mathbf z)\) is
countable for every \(t\in\mathcal T(I)\), for one, equivalently
every, row-choice system \(\mathbf z\).
Let \(\mathcal C\) be the class of objects.

For \(M_0,M_1\in\mathcal C\), declare \(M_0\le_K M_1\) if
\(M_0\preccurlyeq_{\mathrm b}M_1\), and, for every
\(t\in\mathcal T(I^{M_0})\) and every
\(b\in B^{M_1}\setminus B^{M_0}\), the full test for \(t\)
passes when its colors in \(M_1\) are computed using row
representatives in \(M_0\), for \(t\) and all its nonempty
subsequences.  Put \(K=(\mathcal C,\le_K)\).
\end{definition}

Both formulations are well-defined by
Lemma~\ref{aec:representatives}.  The exceptional set in the
object condition is allowed to depend on \(t\); no single
countable set is required to work for all tuples.  In a strong
inclusion, there are no exceptions among the new base-order
points.

\begin{lemma}[Free realization of bit arrays]\label{aec:arrays}
Let \(I\ne\varnothing\), and let \(B\) be an order satisfying
Definition~\ref{aec:basic}(a).  For every family of binary
sequences
\((a_{t,b}\in 2^\omega:t\in\mathcal T(I),\,b\in B)\),
there is a basic structure whose index and base sorts are disjoint
copies of \(I\) and \(B\), with row-choice colors, under these
identifications,
\[
 c_t(b)=[a_{t,b}]_{=^*}.
\]
If \(B\) is countable, every basic structure with base sort
\(B\) is an object.
\end{lemma}

\begin{proof}
Replace \(I\) and \(B\) by disjoint copies and transport the array
indices along these identifications.  Use separate disjoint copies
for the two group sorts as well.
Take tagged, disjoint copies of the appropriate groups for all
the fibers.  Denote the row points by \((t,n,u)\) for
\(u\in[B]^{<\omega}\), and the column points by \((t,b,v)\)
for \(v\in[\omega]^{<\omega}\); the tags distinguish the
carrier types.  Use the regular translation actions and put
\begin{equation}\label{aec:array-realization}
 e((t,n,u),(t,b,v))
 =
 a_{t,b}(n)+\boldsymbol 1_{b\in u}
           +\boldsymbol 1_{n\in v}\pmod2.
\end{equation}
Set bits between different tuples to zero, and equip the
carriers with exactly the specified projections and other
relations.  Translation by symmetric difference gives
\eqref{aec:bit}.  The zero row representatives and a zero
column representative give the sequence \(a_{t,b}\), as
required.  Finally, when \(B\) is countable, each set
\(E_t(\mathbf z)\subseteq B\) is countable automatically.
\end{proof}

\subsection{Isomorphisms, the partial order, and coherence}

\begin{lemma}\label{aec:order-coherence}
The class and strong relation are isomorphism invariant,
\(\le_K\) is a partial order whose comparisons are induced
substructure inclusions, and coherence holds.
\end{lemma}

\begin{proof}
All the basic conditions are preserved by a \(\tau\)-isomorphism,
so the isomorphic image of a basic structure is basic.
Now consider an isomorphism of basic structures.  It
preserves the carrier predicates, every row number \(n\), all
tuple-coordinate projections, the fixed labels of \(G_*\),
the order, supports, actions, and bits.  Transporting a
row-choice system along it thus preserves the color at each
corresponding base point.  Countable bad sets are transported
to countable bad sets.  This proves invariance of the object
condition.  Convexity, fixed \(G_*\), and equality of old
supports are likewise preserved under an isomorphism of a
pair of structures.  Applying the same color calculation at
each new base point proves invariance of the strong relation.
In particular, if \(f:N\cong N'\) and \(M\le_K N\), then
\(f[M]\le_K N'\).

Reflexivity follows because equality is a basic inclusion
and the additional condition has an empty set of new base
points.  Antisymmetry follows from induced substructure
inclusion.  To prove transitivity, suppose
\(M_0\le_K M_1\le_K M_2\).  Basic inclusions compose.
Fix a tuple \(t\) from \(I^{M_0}\), a point
\(b\in B^{M_2}\setminus B^{M_0}\), and row representatives
in \(M_0\) for \(t\) and its nonempty subsequences.
If \(b\in B^{M_1}\setminus B^{M_0}\), choose its column
representatives in \(M_1\).  The bits against these old rows
are unchanged in \(M_2\), by induced substructure, and the
tests hold by \(M_0\le_K M_1\).
If \(b\notin B^{M_1}\), the selected rows in \(M_0\) are
also allowable old representatives for \(M_1\le_K M_2\);
this latter inclusion supplies the tests.  Hence
\(M_0\le_K M_2\).

For coherence, suppose that \(M_0\) is an induced
\(\tau\)-substructure of \(M_1\), and that both are strong
substructures of \(M_2\).  Their \(G_*\) sorts both equal
that of \(M_2\).  Convexity of \(B^{M_0}\) in \(B^{M_2}\)
implies its convexity in \(B^{M_1}\).
For \(u\in G(B)^{M_0}\), its support in \(M_2\) is exactly
its support in \(M_0\); restricting incidence to \(M_1\)
gives the same support there.  Thus
\(M_0\preccurlyeq_{\mathrm b}M_1\).
For a tuple from \(I^{M_0}\) and
\(b\in B^{M_1}\setminus B^{M_0}\), compute colors using
rows in \(M_0\) and columns in \(M_1\).  They are the same
colors computed in \(M_2\), where \(M_0\le_K M_2\) makes
the full test pass.  Therefore \(M_0\le_K M_1\).
\end{proof}

\subsection{Directed unions and smoothness}

\begin{lemma}\label{aec:directed}
Let \((M_i:i\in D)\) be a nonempty set-indexed family of
objects directed under strong inclusion: any two members
have a common strong upper bound within the family.
Its union \(U\) is an object, and \(M_i\le_K U\) for every
\(i\in D\).  If \(N\) is an object with \(M_i\le_K N\)
for all \(i\in D\), then \(U\le_K N\).
\end{lemma}

\begin{proof}
Every finite collection of family members has a common upper
member, by induction from directedness.  In particular,
relations agree on overlapping domains: two members embed
as induced substructures of a common member.  Their union
therefore gives a well-defined relational structure \(U\).

\emph{The basic data in the union.}
The carrier predicates still partition the union domain,
and every finite tuple in \(I^U\) lies in one member.
That member supplies a nonempty row fiber for the tuple at
every \(n<\omega\).  A tuple together with a base point
\(b\) is also contained in a common member, which supplies
the required column fiber.  Projection graphs remain
single-valued and total by the same finite-parameter
argument.

The sorts \(G_*^{M_i}\) are all the same set: for any
two indices, a common upper member has \(G_*\) equal
to both.  Their labels and group structures agree.
Every \(u\in G(B)^U\) originates in a member.  It has
exactly its old support in the union, since a common upper
member with any other proposed incident point cannot add
to that support.  Thus its support is finite.
Conversely, every finite subset of \(B^U\) lies in a
member and is coded there.  Uniqueness of its code can
be checked in a common member containing any two proposed
codes.  The support map in \(U\) is consequently a
bijection onto \([B^U]^{<\omega}\).

All arguments and values needed for a group operation or
an action lie in a common member.  This proves semantic
totality and uniqueness in the union, as well as the
group laws and fiber preservation.  For two row points
in the same union fiber, a common member containing
both supplies their unique translator; the same argument
applies to column points.  Freeness is also checked in
a common member.  The fibers are therefore torsors for
the required groups.  Equation~\eqref{aec:bit} follows
by putting its finitely many arguments in a common
member.  No relations outside the prescribed types
can appear in the union.

The base order is linear, dense, and without endpoints,
since these assertions for finitely many points have
witnesses in a containing member.  More precisely,
if \(a,c\in B^{M_i}\) and \(a<b<c\) in \(B^U\), choose
a member containing \(b\) and then a common upper member
with \(M_i\).  Convexity in that member gives
\(b\in B^{M_i}\).  Hence \(B^{M_i}\) is convex in
\(B^U\).  A bounded interval in the union is exactly
the corresponding interval in any member containing
its endpoints, and so is countable.  These arguments
show that \(U\) is basic and that
\(M_i\preccurlyeq_{\mathrm b}U\) for all \(i\).

\emph{The object condition.}
Fix \(t\in\mathcal T(I^U)\), and choose one member
\(M_i\) containing its entries.  Every row of \(t\)
and of its nonempty subsequences is nonempty already
in \(M_i\).  Choose all their row representatives in
this one member, for every \(n<\omega\).
Over \(B^{M_i}\) their full test has a countable
exception set, by the object condition for \(M_i\).
These same colors can be calculated in \(U\) with
columns from \(M_i\), so the old exception bound
remains valid.

For \(b\in B^U\setminus B^{M_i}\), choose a member
containing \(b\), followed by a common upper member
with \(M_i\).  The strong inclusion of \(M_i\) into
that common member makes the full test pass at \(b\)
using the selected rows.  The same row and column
points compute the same colors in \(U\).  Thus the
test fails only on the original countable set in
\(B^{M_i}\).

The convenient choices just made may depend on \(t\).
To check the definition with a single global
row-choice system in \(U\), compare that system with
these choices for \(t\) and its finitely many
nonempty subsequences.  Lemma~\ref{aec:representatives}
adds at most a countable set of possible differences.
For each \(t\), its bad set is therefore countable
also for the global system.  Hence \(U\) is an object.

\emph{Strongness of the members.}
Fix \(i\), a tuple from \(I^{M_i}\), and
\(b\in B^U\setminus B^{M_i}\).
The same common-upper-member calculation, using
rows from \(M_i\), supplies the full test at \(b\).
Together with the verified basic inclusion this
gives \(M_i\le_K U\).

\emph{A common upper bound.}
Suppose now that every \(M_i\le_K N\).
The union is an induced substructure of \(N\),
because each finite relation instance can be checked
in a member containing its entries.  Its \(G_*\)
sort equals \(G_*^N\), and every union group
element has its original support in \(N\).
For two points of \(B^U\), choose a member
containing both; convexity of that member in \(N\)
forces every intervening point of \(B^N\) into
that member and hence into \(B^U\).  Thus
\(U\preccurlyeq_{\mathrm b}N\).

Fix \(t\in\mathcal T(I^U)\) and
\(b\in B^N\setminus B^U\).
Choose a member \(M_i\) containing \(t\), and
choose all the rows required for \(t\) and its
subsequences inside \(M_i\).
Because \(b\notin B^{M_i}\), the assumption
\(M_i\le_K N\) supplies the full test at \(b\).
Those representatives are also in \(U\) and are
eligible for its strong-inclusion test into \(N\).
Canonicality outside \(B^U\), from
Lemma~\ref{aec:representatives}, makes this one
choice sufficient.  Hence \(U\le_K N\).
\end{proof}

The proof does not put an arbitrary preselected countable
set of union row points into one stage.  It first puts the
finite tuple into a member and then selects all the needed
row representatives there.  This is what permits both the
object-condition argument and smoothness at chains of
countable cofinality.

\subsection{The downward L\"owenheim--Skolem bound}

We first record explicitly the small convex-hull fact
needed for arbitrary requested subsets.

\begin{lemma}\label{aec:convex-hull}
Let \(B\) be a nonempty dense linear order without
endpoints, with countable bounded intervals.  For every
infinite cardinal \(\kappa\) and every \(D\subseteq B\)
of size at most \(\kappa\), there is a nonempty convex
suborder \(B_0\) containing \(D\), of size at most
\(\kappa\), which is dense and has no endpoints.
\end{lemma}

\begin{proof}
If necessary add one point to make \(D\) nonempty.
Set \(T_0=D\).  Given \(T_m\), for each \(x\in T_m\)
choose points \(x^-<x<x^+\) in \(B\), and adjoin them
to form \(T_{m+1}\).  Put
\(T=\bigcup_{m<\omega}T_m\).  Then \(T\ne\varnothing\),
\(|T|\le\kappa\), and \(T\) has no endpoints.
Let
\[
 B_0=\{b\in B:\text{there exist }x,y\in T
                       \text{ with }x\le b\le y\}.
\]
This is the convex hull of \(T\).  It is a union of
at most \(\kappa^2=\kappa\) closed bounded intervals,
each countable after adjoining its endpoints.
Thus \(|B_0|\le\kappa\).
Convexity and density of \(B\) make \(B_0\) dense.
If \(b\in B_0\) lies between \(x,y\in T\), choose
points of \(T\) below \(x\) and above \(y\).
They give points of \(B_0\) on both sides of \(b\),
so \(B_0\) has no endpoints.
\end{proof}

\begin{lemma}\label{aec:ls}
For every \(M\in\mathcal C\) and every
\(A\subseteq |M|\), there is \(N\le_K M\) such that
\(A\subseteq |N|\) and
\(\|N\|\le |A|+\aleph_0\).
\end{lemma}

\begin{proof}
Put \(\kappa=|A|+\aleph_0\).
Choose a nonempty \(I_0\subseteq I^M\) of size
at most \(\kappa\) containing \(A\cap I^M\)
and every tuple-coordinate projection of every
row or column element of \(A\).
There are at most \(\kappa\) tuples in
\(\mathcal T(I_0)\).  For each of them, and for
each \(n<\omega\), choose a representative
\(z_{t,n}\in Z^M_{t,n}\).
When computing the bad sets \(E_t\) for
\(t\in\mathcal T(I_0)\), always use this same
system of choices for all its subtuples.
The object condition in \(M\), or an extension
of these choices to a full row-choice system
in \(M\), shows that each such \(E_t\) is
countable.

Collect a subset \(D\subseteq B^M\) consisting of:
\begin{enumerate}[label=\textup{(\roman*)}]
\item all points of \(A\cap B^M\), and all
base-point projections of column elements of \(A\);
\item \(\supp(u)\) for every
\(u\in A\cap G(B)^M\);
\item for every requested row point
\(z\in A\cap Z^M_{t,n}\), the support of the
unique \(u\in G(B)^M\) satisfying
\(z=z_{t,n}+u\);
\item all the sets \(E_t\), for
\(t\in\mathcal T(I_0)\).
\end{enumerate}
In (ii)--(iv) the listed sets are included by
taking their union.  The first three clauses
contribute at most \(\kappa\) points, since
supports and projections are finite for each
requested element.  The last clause contributes
at most \(\kappa\cdot\aleph_0=\kappa\)
points.  Therefore \(|D|\le\kappa\).

Apply Lemma~\ref{aec:convex-hull} to obtain a
nonempty convex dense \(B_0\subseteq B^M\)
without endpoints, containing \(D\), of size
at most \(\kappa\).
Let
\[
 H=\{u\in G(B)^M:\supp^M(u)\subseteq B_0\}.
\]
This is exactly the subgroup coding
\([B_0]^{<\omega}\).
Define \(N\) by retaining \(I_0,B_0,H\), and
the whole group \(G_*^M\), together with the
fibers
\begin{align*}
 Z^N_{t,n}&=\{z_{t,n}+u:u\in H\}
       &&(t\in\mathcal T(I_0),\ n<\omega),\\
 W^N_{t,b}&=W^M_{t,b}
       &&(t\in\mathcal T(I_0),\ b\in B_0),
\end{align*}
and all relations restricted from \(M\).

Each displayed row orbit is a nonempty
\(H\)-torsor.  Each retained column is the
whole old \(G_*^M\)-torsor and remains one.
The projections land in the retained sorts,
the group operations and actions are total
there, and Equation~\eqref{aec:bit} restricts
to these groups and fibers.  Thus \(N\) is
basic and \(N\preccurlyeq_{\mathrm b}M\).

The choices \(z_{t,n}\) are present in \(N\):
they are obtained by the zero translator.
At every \(b\in B_0\), they compute precisely
the same colors in \(N\) and \(M\), using the
retained whole column fibers.  The bad set in
\(N\) for any \(t\in\mathcal T(I_0)\) is
therefore \(E_t\cap B_0\), which is countable.
Hence \(N\) is an object.
For \(b\in B^M\setminus B_0\), all these
full tests pass because \(E_t\subseteq D
\subseteq B_0\).  The same chosen rows are
eligible old representatives from \(N\).
This proves the strong condition, so
\(N\le_K M\).

It remains to verify containment of the
requested data.  The choices of \(I_0\)
and \(B_0\) retain requested index and base
points and all needed projections.  A
requested group element belongs to \(H\)
by clause (ii).  A requested row point
belongs to its retained orbit by clause
(iii).  A requested column point belongs
to a retained whole column fiber by
clause (i) and the choice of \(I_0\).
All requested elements of \(G_*^M\)
are present because that entire group was
retained.  These are all possible carrier
types, so \(A\subseteq |N|\).
Finally Lemma~\ref{aec:size} gives
\[
 \|N\|=\max(|I_0|,|B_0|,\aleph_0)
       \le\kappa,
\]
as required.
\end{proof}

\begin{theorem}\label{aec:class}
The pair \(K=(\mathcal C,\le_K)\) of
Definition~\ref{aec:objects} is an abstract
elementary class in a fixed countable finitary
relational language, with
\[
 \LS(K)=\aleph_0.
\]
It has the union and common-upper-bound
properties for every nonempty set-indexed
directed family of strong substructures.
\end{theorem}

\begin{proof}
The language and its semantic structures were
specified after Definition~\ref{aec:basic}.
Lemma~\ref{aec:order-coherence} proves
isomorphism invariance, substructure, the
partial-order properties, and coherence.
Lemma~\ref{aec:directed} proves the directed
union and smoothness assertions, in
particular the required chain axioms.
Lemma~\ref{aec:ls} proves the downward
bound \(\aleph_0\).  Since the language is
countable and \(\LS(K)\) is by definition
infinite, its least possible value is
\(\aleph_0\), which is therefore the value
attained here.
\end{proof}
